\section{Simulation in COMSOL Multiphysics}

\subsection{Purpose of the COMSOL Simulation}

The purpose of the COMSOL Multiphysics simulation was to generate a controlled synthetic dataset for the study of thermoelastic wave propagation in geological media. Datasets that would be suitable for the needs of model training were not found in the public domain. Therefore, it was decided to generate synthetic datasets in third-party programs.

The required dataset had to contain not only temperature values, but also displacement, velocity, stress, strain, material parameters, spatial coordinates, and mesh information. These quantities are difficult to obtain experimentally at every spatial point and every time step. In contrast, a numerical multiphysics model allows all required physical fields to be extracted consistently from the same simulation.

in the process of searching for programs, the programs Ansys, Nvidia Omniverse, Comsol Multiphysics were found. Ansys was considered as a powerful finite element analysis platform, but the construction of a repeatable data-generation pipeline required more computational resources and a more complex workflow. Nvidia Omniverse was also considered for visualization purposes; however, the main goal of this work was not only to visualize the process, but also to export numerical fields for further processing. Comsol Multiphysics was selected because it supports coupled thermo-mechanical analysis, finite element discretization, time-dependent studies, material assignment, boundary condition configuration, visualization, and structured export of simulation data. Also, when configuring Comsol, it turned out to be more accessible in the absence of computing capabilities.

\begin{figure}[H]
\centering
\includegraphics[width=0.85\textwidth]{chapters/chapter4/comsol_interface.png}
\caption{COMSOL Multiphysics interface used for creating the thermoelastic simulation.}
\label{fig:comsol_interface}
\end{figure}

\subsection{Geometry of the Thermoelastic Experiment}

The simulated experiment represents a two-dimensional geological medium with a heated metallic rod inserted into the rock domain. The rock is treated as the main geological body, while the rod acts as a localized thermal excitation source. The purpose of this setup is to create a non-uniform temperature field, which then produces thermal expansion, mechanical stress, displacement, and wave-like thermoelastic response in the surrounding medium.

The computational domain is denoted as

\begin{equation}
\Omega =
\Omega_{\mathrm{rock}}
\cup
\Omega_{\mathrm{rod}} .
\end{equation}

Here, \(\Omega\) is the full computational domain, \(\Omega_{\mathrm{rock}}\) is the rock region, and \(\Omega_{\mathrm{rod}}\) is the heated rod region. These two domains are assigned different thermal and mechanical material properties in COMSOL.

The initial temperature of the rock is denoted by \(T_0\), while the initial temperature of the rod is denoted by \(T_{\mathrm{rod}}\). The main thermal excitation condition is

\begin{equation}
T_{\mathrm{rod}} > T_0 .
\end{equation}

In this expression, \(T_{\mathrm{rod}}\) is the initial temperature of the metallic rod, and \(T_0\) is the initial temperature of the surrounding rock. This temperature contrast creates heat conduction from the rod into the geological medium. Because the temperature distribution becomes spatially non-uniform, different parts of the rock expand by different amounts. This produces strain, stress, and displacement fields.

\subsection{COMSOL Model Structure}

The simulation was created as a coupled thermo-mechanical model. The structure of the COMSOL project included global parameters, material definitions, geometry, physical interfaces, multiphysics coupling, mesh generation, and a time-dependent study.

The main physics interfaces used in the model were:

The COMSOL simulation was constructed using several physics interfaces and model components. Each of them had a specific role in describing the coupled thermoelastic behavior of the geological medium.

Solid Mechanics.
The Solid Mechanics interface was used to calculate the mechanical behavior of the rock domain. It calculated displacement, velocity, strain, stress and propagation of mechanical disturbances due to thermal expansion. This interface was necessary because heating of the rod creates deformation in the surrounding geological medium.

Linear Elastic Material.
The Linear Elastic Material model defined the mechanical behavior of the rock. It described the medium as an elastic solid using material parameters such as Young's modulus, Poisson's ratio, density, and the thermal expansion coefficient. These parameters determined how strongly the rock deformed under thermoelastic loading.

Free Boundary.
The Free Boundary condition was applied to selected external boundaries of the geometry. It allowed these boundaries to deform without externally applied mechanical traction. This condition was used to reproduce a natural mechanical response of the medium instead of artificially fixing all outer surfaces.

Heat Transfer in Solids.
The Heat Transfer in Solids interface was used to model transient heat conduction in the simulation domain. It calculated how heat from the heated rod diffused into the surrounding rock with time. This module provided the temperature field needed for the thermoelastic coupling.

Thermal Insulation.
Thermal Insulation condition was imposed on some external boundaries. It prevented heat flux through these surfaces and reduced the heat exchange with the outside environment. This made it possible to control the thermal behavior of the model and focus on heat propagation inside the simulated rock domain.

Thermal Expansion.
The coupling between the thermal and mechanical parts of the simulation was made using the Thermal Expansion feature. It converted temperature changes into thermal strain. The temperature field obtained from the heat transfer analysis then led to mechanical deformation in the solid mechanics model.

Temperature Coupling.
The Temperature Coupling feature transferred the computed temperature field from the Heat Transfer in Solids interface to the Solid Mechanics interface. This ensured that the mechanical response of the rock was calculated using the actual transient temperature distribution inside the model.

Time Dependent Study.
The Time Dependent Study was used to solve the coupled thermoelastic problem over a sequence of time steps. It generated time-dependent temperature, displacement, stress and strain fields. These r sults were later exported and used for synthetic dataset generation.

This model structure is important because the final dataset is not generated from independent physical fields. Temperature, displacement, stress, and strain are obtained from the same coupled solution, which makes the exported data physically consistent.

\begin{figure}[H]
\centering
\includegraphics[width=0.85\textwidth]{chapters/chapter4/comsol_model_tree.png}
\caption{COMSOL model tree with Solid Mechanics, Heat Transfer in Solids, multiphysics coupling, mesh, and time-dependent study.}
\label{fig:comsol_model_tree}
\end{figure}
The COMSOL experiment was designed as a controlled thermoelastic scenario in which a highly heated rod interacts with a geological medium. The main goal of this setup was to create a localized thermal disturbance inside the rock and to observe how this disturbance produces heat conduction, thermal expansion, stress formation, displacement, and wave-like mechanical response.

The main computational domain was represented by a square rock block with a side length of \(1~\mathrm{m}\). Therefore, the model describes a \(1~\mathrm{m} \times 1~\mathrm{m}\) cross-section of a geological medium. This block was used as the main rock domain in which the thermoelastic response was studied.

A cylindrical rod was placed inside the rock block and used as the localized heat source. The rod had a radius of \(0.04~\mathrm{m}\) and a length of \(0.5~\mathrm{m}\). Since the size of the main block was \(1~\mathrm{m}\), the rod occupied half of the characteristic length of the domain. This configuration made it possible to heat only a limited internal region of the rock instead of applying a uniform temperature load to the whole medium.

At the beginning of the simulation, the rock was initialized at room temperature:

\begin{equation}
T_0 = 293.15~\mathrm{K}.
\end{equation}

The reference temperature for thermal expansion was also set to the same value:

\begin{equation}
T_{\mathrm{ref}} = 293.15~\mathrm{K}.
\end{equation}

This means that the rock was considered thermally undeformed at the initial state. In other words, before the heated rod affected the medium, no thermal strain was assumed to exist in the rock.

The rod temperature was set much higher than the initial rock temperature:

\begin{equation}
T_{\mathrm{rod},0} = 1500~\mathrm{K}.
\end{equation}

As a result, a strong temperature contrast appeared between the rod and the surrounding geological medium. The temperature difference between the heated rod and the initial rock state was

\begin{equation}
\Delta T = 1500 - 293.15 = 1206.85~\mathrm{K}.
\end{equation}

This large temperature contrast created an intensive thermal gradient near the rod--rock contact region. Heat started to propagate from the rod into the surrounding medium. Since the temperature field was not uniform, different parts of the rock expanded by different amounts. This non-uniform expansion generated thermal strain and stress in the material.

The physical logic of the experiment can be described as follows. First, the heated rod transfers heat to the surrounding rock. Then, the temperature increase causes local thermal expansion. Because the heated region is surrounded by colder material, this expansion is partially constrained. As a consequence, thermal stresses appear around the rod. These stresses lead to displacement of the rock material and form the basis of the thermoelastic wave-like response observed in the simulation.

The main parameters of the experiment are presented in Table~\ref{tab:heated_rod_experiment_parameters}.

\begin{table}[H]
\centering
\caption{Main parameters of the heated rod thermoelastic experiment.}
\label{tab:heated_rod_experiment_parameters}
\begin{tabular}{p{0.35\textwidth} p{0.25\textwidth} p{0.25\textwidth}}
\hline
Parameter & Symbol & Value \\
\hline
Main rock block size & \(L\) & \(1~\mathrm{m}\) \\
Rod radius & \(r_{\mathrm{rod}}\) & \(0.04~\mathrm{m}\) \\
Rod length & \(L_{\mathrm{rod}}\) & \(0.5~\mathrm{m}\) \\
Initial rock temperature & \(T_0\) & \(293.15~\mathrm{K}\) \\
Initial rod temperature & \(T_{\mathrm{rod},0}\) & \(1500~\mathrm{K}\) \\
Initial temperature contrast & \(\Delta T\) & \(1206.85~\mathrm{K}\) \\
\hline
\end{tabular}
\end{table}

\subsection{Thermoelastic Physical Model}

The thermal part of the simulation is based on transient heat conduction. The governing equation can be written as

\begin{equation}
\rho C_p \frac{\partial T}{\partial t}
=
\nabla \cdot \left( k \nabla T \right) + Q .
\end{equation}

Here, \(T\) is temperature, \(t\) is time, \(\rho\) is material density, \(C_p\) is heat capacity at constant pressure, \(k\) is thermal conductivity, and \(Q\) is a volumetric heat source. The term on the left describes thermal energy accumulation, while the term \(\nabla \cdot (k \nabla T)\) describes heat conduction inside the solid material.

The mechanical response is computed by the Solid Mechanics interface. The displacement field is denoted as \(\mathbf{u}\). In the two-dimensional formulation, the main in-plane displacement components are \(u\) and \(v\). The COMSOL export also stores additional available components such as \(w\) and \(w_t\), which are retained during data processing to preserve the complete solver output.

Thermoelastic coupling is introduced through thermal strain:

\begin{equation}
\boldsymbol{\varepsilon}_{\mathrm{th}}
=
\alpha
\left(
T - T_{\mathrm{ref}}
\right)
\mathbf{I}.
\end{equation}

In this formula, \(\boldsymbol{\varepsilon}_{\mathrm{th}}\) is thermal strain, \(\alpha\) is the coefficient of thermal expansion, \(T\) is the current temperature, \(T_{\mathrm{ref}}\) is the reference temperature, and \(\mathbf{I}\) is the identity tensor. The formula means that a temperature increase produces expansion of the material. This thermal expansion changes the stress and displacement fields.

The dynamic mechanical response is described in COMSOL by the balance of linear momentum:

\begin{equation}
\rho_s
\frac{\partial^2 \mathbf{u}}{\partial t^2}
=
\nabla \cdot \boldsymbol{\sigma}
+
\mathbf{f}.
\end{equation}

Here, \(\rho_s\) is the structural density, \(\mathbf{u}\) is the displacement vector, \({\sigma}\) is the stress tensor, and \(\mathbf{f}\) is the body force vector. This equation states that stress gradients and external forces produce acceleration of the medium. In the present experiment, the main source of mechanical response is not an external impact, but thermal expansion caused by the heated rod.

The coupling chain of the simulation can therefore be summarized as follows:

\begin{equation}
T
\rightarrow
\boldsymbol{\varepsilon}_{\mathrm{th}}
\rightarrow
\boldsymbol{\sigma}
\rightarrow
\mathbf{u}.
\end{equation}

This expression shows the main physical logic of the experiment: the temperature field creates thermal strain, thermal strain modifies the stress field, and the stress field produces displacement and wave-like mechanical response.

\subsection{Initial and Boundary Conditions}

The initial conditions define the temperature contrast between the heated rod and the surrounding rock. The rock domain is initialized with temperature \(T_0\), while the rod domain is initialized with a higher temperature \(T_{\mathrm{rod}}\):

\begin{equation}
T(\mathbf{x},0) = T_0,
\qquad
\mathbf{x} \in \Omega_{\mathrm{rock}},
\end{equation}

\begin{equation}
T(\mathbf{x},0) = T_{\mathrm{rod}},
\qquad
\mathbf{x} \in \Omega_{\mathrm{rod}}.
\end{equation}

Here, \(\mathbf{x}\) denotes a spatial point in the computational domain, \(T(\mathbf{x},0)\) is the initial temperature at this point, \(T_0\) is the initial rock temperature, and \(T_{\mathrm{rod}}\) is the initial rod temperature.

For the external thermal boundaries, thermal insulation was applied. This condition can be written as

\begin{equation}
-\mathbf{n} \cdot \mathbf{q} = 0 .
\end{equation}

In this expression, \(\mathbf{n}\) is the outward normal vector to the boundary, and \(\mathbf{q}\) is the heat flux vector. The condition means that no heat leaves the model through the insulated boundary. This allows the simulation to focus on heat propagation inside the rock domain.

For the mechanical part, free boundary conditions were used on selected boundaries. A free mechanical boundary means that no external traction is applied:

\begin{equation}
\boldsymbol{\sigma}\mathbf{n} = 0 .
\end{equation}

Here, \({\sigma}\) is the stress tensor and \(\mathbf{n}\) is the outward boundary normal. This condition allows the boundary to deform naturally under the internal thermoelastic response.

\subsection{Material Parameters}

The material parameters define how heat propagates through the medium and how the medium mechanically responds to thermal loading. The rock and the rod are assigned different properties. The most important parameters are listed in Table~\ref{tab:comsol_material_parameters}.

\begin{table}[H]
\centering
\caption{Main material parameters used in the thermoelastic simulation.}
\label{tab:comsol_material_parameters}
\begin{tabular}{p{0.26\textwidth} p{0.16\textwidth} p{0.48\textwidth}}
\hline
Parameter & Symbol & Physical meaning \\
\hline
Density & \(\rho\), \(\rho_s\) & Controls thermal storage and mechanical inertia. In COMSOL, density can be exported separately for heat transfer and solid mechanics. \\
\hline
Thermal conductivity & \(k\) & Controls the rate of heat diffusion through the material. \\
\hline
Heat capacity & \(C_p\) & Defines how much energy is required to change the material temperature. \\
\hline
Young's modulus & \(E\) & Defines the stiffness of the solid material. \\
\hline
Poisson's ratio & \(\nu\) & Describes lateral deformation under mechanical loading. \\
\hline
Thermal expansion coefficient & \(\alpha\) & Couples temperature change with mechanical strain. \\
\hline
\end{tabular}
\end{table}

These parameters are essential for dataset generation because changing them changes both the temperature field and the mechanical response. For example, a material with high thermal conductivity transfers heat faster, while a material with a high Young's modulus produces a different stress and displacement response under the same temperature loading.

\subsection{Mesh Generation and Time-Dependent Study}

The coupled thermoelastic problem was solved using the finite element method. In COMSOL, the continuous geometry is divided into a finite element mesh, and the solution is computed over this discretized domain. The mesh is especially important near the heated rod because this region contains the strongest temperature gradients and stress concentrations.

The simulation was performed as a time-dependent study. The simulated time interval is represented as

\begin{equation}
t_n = n \Delta t,
\qquad
n = 0,1,\ldots,N_t .
\end{equation}

Here, \(t_n\) is the time value at step \(n\), \(\Delta t\) is the time step, and \(N_t\) is the total number of simulated time steps. At every time step, COMSOL computes the temperature field, displacement field, velocity field, stress tensor components, strain components, and material-related quantities.

The result of the simulation is therefore not a single static solution. It is a time-dependent sequence of physical states that describes how heat and mechanical disturbances evolve inside the geological medium.

\subsection{COMSOL Data Export and Dataset Processing}

The raw dataset was exported from COMSOL as a set of CSV files. Each export used the dataset \textit{Study 1/Solution 1 (sol1)} with time selection set to \textit{All}. Therefore, the exported tables contain the full time-dependent solution rather than a single selected time step. The mesh was exported separately as a COMSOL Multiphysics text file with the extension \texttt{.mphtxt}. This file preserves the spatial discretization of the numerical model and can later be used to reconstruct the mesh-based graph representation.

The exported files were stored in the raw dataset directory:
\begin{equation}
\texttt{datasets/<dataset\_id>/raw/}.
\end{equation}
The raw directory contains CSV files with physical fields, the exported mesh file, and a scenario configuration file. An example of the raw dataset structure is shown in Figure~\ref{fig:raw_dataset_folder}.

\begin{figure}[H]
    \centering
    \includegraphics[width=0.85\textwidth]{chapters/chapter4/raw_dataset_folder.png}
    \caption{Raw dataset structure exported from COMSOL, including physical field CSV files, mesh data, and scenario configuration.}
    \label{fig:raw_dataset_folder}
\end{figure}

The exported COMSOL groups are summarized in Table~\ref{tab:comsol_export_groups}. The files are separated according to their physical meaning: temperature-related quantities, displacement and velocity fields, material parameters, stress components, strain components, and mesh information.

\begin{table}[H]
\centering
\caption{Exported COMSOL data groups used for dataset construction.}
\label{tab:comsol_export_groups}
\begin{tabular}{p{0.24\textwidth} p{0.34\textwidth} p{0.32\textwidth}}
\hline
\textbf{Export label} & \textbf{Exported expressions} & \textbf{Description} \\
\hline
\texttt{data\_temperature} & 
\texttt{T}, \texttt{x}, \texttt{y}, \texttt{z}, \texttt{ht.k\_iso}, \texttt{ht.rho}, \texttt{ht.Cp} &
Temperature field, spatial coordinates, and thermal material properties. \\
\hline
\texttt{data\_displacement} &
\texttt{u}, \texttt{v}, \texttt{w}, \texttt{ut}, \texttt{vt}, \texttt{wt} &
Displacement and structural velocity components. \\
\hline
\texttt{data\_materials} &
\texttt{solid.E}, \texttt{solid.nu}, \texttt{solid.rho}, \texttt{te1.alpha\_iso} &
Elastic and thermoelastic material properties. \\
\hline
\texttt{data\_stress\_1} &
\texttt{solid.mises}, \texttt{solid.sx}, \texttt{solid.sy}, \texttt{solid.sz} &
Von Mises stress and normal stress components. \\
\hline
\texttt{data\_stress\_2} &
\texttt{solid.sxy}, \texttt{solid.syz}, \texttt{solid.sxz} &
Shear stress components. \\
\hline
\texttt{data\_strain} &
\texttt{solid.eX}, \texttt{solid.eY}, \texttt{solid.eZ}, \texttt{solid.eXY}, \texttt{solid.eYZ}, \texttt{solid.eXZ} &
Strain tensor components exported from Solid Mechanics. \\
\hline
\texttt{mesh\_rod\_experiment} &
\texttt{.mphtxt} mesh file &
Finite element mesh of the rod experiment. \\
\hline
\end{tabular}
\end{table}

Each CSV file contains COMSOL metadata rows and a numerical table. The metadata rows begin with the symbol \(\texttt{\%}\) and describe the model name, COMSOL version, date of export, spatial dimension, number of nodes, and exported expressions. The actual data table starts from the row containing the spatial coordinates
\begin{equation}
X,\quad Y,\quad Z.
\end{equation}
The remaining columns contain physical variables written in the COMSOL format
\begin{equation}
\texttt{field\_name (unit) @ t=value}.
\end{equation}
This naming convention stores three important pieces of information in each column name: the physical variable, its unit of measurement, and the simulation time. An example of such a raw COMSOL export is shown in Figure~\ref{fig:raw_comsol_csv_example}.

\begin{figure}[H]
    \centering
    \includegraphics[width=0.95\textwidth]{chapters/chapter4/example_of_dataset.png}
    \caption{Example of a raw COMSOL CSV export containing metadata rows, spatial coordinates, physical field names, units, and time-dependent columns.}
    \label{fig:raw_comsol_csv_example}
\end{figure}

After export, the raw numerical files were converted into a unified machine-learning dataset. This step was necessary because COMSOL tables are solver-oriented: physical fields are stored in wide CSV format, metadata rows are mixed with numerical values, and different quantities are split across separate files. Neural surrogate models require a cleaner representation with aligned spatial points, synchronized time steps, normalized physical channels, and a clear separation between dynamic and static information.

The preprocessing stage was designed as a common data preparation step for all neural models considered in this work. Graph-based models can use the mesh connectivity and edge attributes, operator-learning models can use the spatio-temporal field tensor, and coordinate-based models can sample points in space and time. This allows different architectures to be compared using the same physical dataset rather than separate preprocessing assumptions.

During CSV cleaning, COMSOL service rows were skipped, the true header row was detected, and only supported physical exports were parsed. Non-numerical entries were converted where possible, while invalid or missing values were handled before tensor construction. This was required because the exported variables have different units and numerical scales: temperature is measured in kelvin, displacement in meters, stress in pascals, strain is dimensionless, and material parameters may have substantially larger magnitudes.

A separate coordinate-alignment step was required because different COMSOL exports may contain different point sets or the same points in a different order. Therefore, the files cannot be merged reliably by row index. A reference coordinate array was selected, and the remaining fields were mapped to this reference using nearest-neighbor matching in coordinate space:
\begin{equation}
j(i)
=
\arg\min_j
\left\|
\mathbf{r}^{\mathrm{ref}}_i
-
\mathbf{r}^{\mathrm{src}}_j
\right\|_2 .
\end{equation}
Here, \(\mathbf{r}^{\mathrm{ref}}_i\) is the \(i\)-th coordinate in the reference node set, and \(\mathbf{r}^{\mathrm{src}}_j\) is a coordinate from another exported file. This step produced a single aligned node set shared by all physical fields.

After alignment, the data were separated into dynamic and static features. Dynamic features describe the time-dependent state of the thermoelastic process. These include temperature, displacement components, velocity components, stress components, and strain components. Static features describe quantities that do not change during one simulation, such as spatial coordinates, material parameters, source-region indicators, distance to the source, and scenario-level parameters. Material properties are not treated as prediction targets; they are used as conditioning information for the models.

The dynamic part of the dataset was assembled into a spatio-temporal tensor
\begin{equation}
\mathbf{S}
\in
\mathbb{R}^{N_t \times N_p \times F_d},
\end{equation}
where \(N_t\) is the number of time steps, \(N_p\) is the number of spatial points or mesh nodes, and \(F_d\) is the number of dynamic physical channels. For one time step \(n\), the state of the simulated medium is represented as
\begin{equation}
\mathbf{S}_n
\in
\mathbb{R}^{N_p \times F_d}.
\end{equation}
The static part was stored as
\begin{equation}
\mathbf{A}
\in
\mathbb{R}^{N_p \times F_s},
\end{equation}
where \(F_s\) is the number of static node or scenario features.

Normalization was applied separately to dynamic and static features. For each feature channel, the mean and standard deviation were computed and saved for later denormalization. The normalized value was calculated as
\begin{equation}
\widetilde{x}
=
\frac{x-\mu}{\sigma+\varepsilon},
\end{equation}
where \(x\) is the original feature value, \(\mu\) is the mean value of the corresponding channel, \(\sigma\) is its standard deviation, and \(\varepsilon\) is a small constant used for numerical stability. Separate normalization statistics are important because dynamic physical fields and static material parameters have different distributions and units.

For graph-based models, the aligned coordinates were additionally converted into a graph representation. If the exported \(\texttt{.mphtxt}\) file was available and compatible with the aligned node set, mesh connectivity was used to construct graph edges. If the mesh connectivity could not be used, a nearest-neighbor graph was constructed from the node coordinates. For an edge from node \(i\) to node \(j\), the edge attribute vector was defined as
\begin{equation}
\mathbf{e}_{ij}
=
\left[
x_j-x_i,\,
y_j-y_i,\,
z_j-z_i,\,
\left\|
\mathbf{r}_j-\mathbf{r}_i
\right\|_2
\right].
\end{equation}
The first three components describe the relative spatial offset between neighboring nodes, while the last component stores the Euclidean distance.

Training pairs were generated from the processed time sequence. Depending on the model type, the processed dataset can be used as full-field states, coordinate-time samples, temporal sequences, or graph trajectories. For one-step transition learning, the target can be represented as the state increment
\begin{equation}
\Delta \mathbf{S}_n
=
\mathbf{S}_{n+1}
-
\mathbf{S}_n .
\end{equation}
This delta representation is useful for transient thermoelastic dynamics because the model learns the local change of the physical state between consecutive time steps rather than predicting the complete field independently at every step.

The output of the preprocessing stage was stored in
\begin{equation}
\texttt{datasets/<dataset\_id>/processed/}.
\end{equation}
An example of the processed dataset structure is shown in Figure~\ref{fig:processed_dataset_folder}.

\begin{figure}[H]
    \centering
    \includegraphics[width=0.85\textwidth]{chapters/chapter4/processed_dataset_folder.png}
    \caption{Processed dataset artifacts generated after parsing, coordinate alignment, normalization, graph construction, and trajectory preparation.}
    \label{fig:processed_dataset_folder}
\end{figure}

The most important processed artifacts are summarized in Table~\ref{tab:processed_dataset_artifacts}.

\begin{table}[H]
\centering
\caption{Processed dataset artifacts produced from COMSOL exports.}
\label{tab:processed_dataset_artifacts}
\begin{tabular}{p{0.32\textwidth} p{0.58\textwidth}}
\hline
\textbf{Artifact} & \textbf{Purpose} \\
\hline
\texttt{coords.npy} &
Aligned spatial coordinates of the computational points or mesh nodes. \\
\hline
\texttt{dynamic.npy} &
Time-dependent physical fields stored as a spatio-temporal tensor. \\
\hline
\texttt{static.npy} &
Static node, material, and scenario features. \\
\hline
\texttt{time.npy} &
Simulation time values corresponding to the dynamic states. \\
\hline
\texttt{normalization.json} &
Normalization statistics required for scaling and denormalization. \\
\hline
\texttt{metadata.json} &
Metadata for the dataset: feature names, dimensions, time-step information, processing. \\
\hline
\texttt{graph.pt} &
Graph representation used by graph-based models: edge index, edge attributes, coordinates, and static node features. \\
\hline
\texttt{trajectories.pt} &
Prepared transition samples or trajectory fragments for models that learn temporal state evolution. \\
\hline
\texttt{preview.csv} &
Diagnostic summary used to inspect feature ranges and verify the correctness of preprocessing. \\
\hline
\end{tabular}
\end{table}

Thus, the preprocessing stage converts heterogeneous COMSOL outputs into a common numerical dataset that can be reused by different neural surrogate models. This design is important for fair comparison: the models differ in architecture and input representation, but they are trained and evaluated from the same cleaned, aligned, and normalized thermoelastic simulation data.

\subsection{Visualization and Validation of the Simulation}

Visualization was used to check whether the simulation behaved physically correctly before using the exported data. The most important checks were the geometry, mesh quality, temperature propagation, and stress distribution.

The temperature field was inspected to verify that heat propagates from the heated rod into the surrounding rock. This check is important because an incorrect material assignment, wrong thermal contact, or incorrect boundary condition can prevent heat from entering the rock domain.

\begin{figure}[H]
\centering
\includegraphics[width=0.85\textwidth]{chapters/chapter4/temperature.png}
\caption{Temperature distribution in the rock caused by the heated rod.}
\label{fig:comsol_temperature}
\end{figure}

The stress field was inspected to verify that the non-uniform temperature distribution creates a mechanical response. Regions close to the rod are expected to show stronger stress because they experience the highest thermal gradients and strongest local expansion.

\begin{figure}[H]
\centering
\includegraphics[width=0.85\textwidth]{chapters/chapter4/stress_1.png}
\caption{Stress distribution at the beginning of the thermoelastic response.}
\label{fig:comsol_stress_beginning}
\end{figure}

\begin{figure}[H]
\centering
\includegraphics[width=0.85\textwidth]{chapters/chapter4/stress_2.png}
\caption{Stress distribution during the intermediate stage of the simulation.}
\label{fig:comsol_stress_middle}
\end{figure}

\begin{figure}[H]
\centering
\includegraphics[width=0.85\textwidth]{chapters/chapter4/stress_3.png}
\caption{Stress distribution at the later stage of the simulation.}
\label{fig:comsol_stress_end}
\end{figure}

These visualizations were used as a qualitative validation step. If the temperature field did not propagate from the rod, or if the stress field remained physically unrealistic, the simulation configuration had to be corrected before exporting the dataset.

\subsection{Limitations of the Simulation Setup}

The generated dataset depends on the assumptions used in the COMSOL model. The simulation considers a controlled two-dimensional thermoelastic setup, which is suitable for generating consistent training data but does not fully reproduce all complexities of real geological formations. Real rocks may contain cracks, pores, layers, anisotropy, and heterogeneous material properties.

The exported fields' accuracy is also affected by the mesh resolution. A coarse mesh may not be able to resolve sharp gradients near the heated rod. A very fine mesh is computationally expensive and can result in large file size. The time step has a similar effect: smaller time steps provide a more detailed temporal description, but they produce larger CSV files and require more processing time.

Another limitation is that the dataset is synthetic. It is physically consistent within the assumptions of the numerical model, but it is not a direct measurement from a real geological experiment. Nevertheless, the synthetic approach is appropriate for this work because it provides complete access to temperature, displacement, velocity, stress, strain, material parameters, coordinates, and mesh information at all simulated time steps.

The thermoelastic experiment was mainly designed and realized through COMSOL Multiphysics, which was used to solve the coupled physical problem, visualize the propagation process and export the structured dataset for further computational analysis.